\documentclass[11pt,a4paper]{article}
\usepackage[T1]{fontenc}\usepackage[utf8]{inputenc}\usepackage{lmodern}
\usepackage[margin=23mm,headheight=15pt]{geometry}
\usepackage{amsmath,amssymb,amsthm,mathtools,booktabs,longtable,array,microtype,xcolor}
\usepackage{hyperref,fancyhdr,listings,enumitem}
\definecolor{ink}{HTML}{163348}\definecolor{accent}{HTML}{147D83}
\hypersetup{colorlinks=true,linkcolor=ink,urlcolor=accent,citecolor=accent,pdfauthor={Sébastien Palcoux}}
\pagestyle{fancy}\fancyhf{}\fancyhead[L]{\sffamily\small FUSION RING CENSUS}
\fancyhead[R]{\sffamily\small Rank \Rank}\fancyfoot[L]{\sffamily\footnotesize Computational companion \textbullet\ September 2026}
\fancyfoot[R]{\sffamily\small\thepage}\renewcommand{\headrulewidth}{0.4pt}
\newtheorem{theorem}{Theorem}[section]\newtheorem{proposition}[theorem]{Proposition}
\newtheorem{lemma}[theorem]{Lemma}\newtheorem{corollary}[theorem]{Corollary}
\theoremstyle{definition}\newtheorem{definition}[theorem]{Definition}\newtheorem{remark}[theorem]{Remark}
\newcommand{\ZZ}{\mathbb Z}\newcommand{\CC}{\mathbb C}\newcommand{\Q}{\mathbb Q}
\newcommand{\NN}{\mathbb Z_{\geq0}}\newcommand{\FPdim}{\operatorname{FPdim}}
\setlength{\parindent}{0pt}\setlength{\parskip}{0.45em}\setlength{\emergencystretch}{3em}
\lstset{basicstyle=\ttfamily\footnotesize,breaklines=true,columns=fullflexible,frame=single,rulecolor=\color{accent},aboveskip=1em,belowskip=1em}
\setcounter{tocdepth}{2}
\newcommand{\Rank}{3}
\begin{document}
\thispagestyle{empty}{\sffamily\small\color{accent} EXACT ENUMERATION / DATA / VERIFICATION}\par\vspace{1cm}
{\sffamily\Huge\bfseries\color{ink} Rank-3 fusion rings\par}
\vspace{0.35cm}
{\sffamily\Large A reproducible census through multiplicity 1000\par}
\vspace{0.65cm}{\large Sébastien Palcoux}\par{\small BIMSA}\par\vspace{0.35cm}{\small Computational companion, 21 September 2026}\par\vspace{0.8cm}
\begin{abstract}
We explain the exact enumeration of all based fusion rings of rank 3 through multiplicity 1000, comprising 460,353 isomorphism classes. The account separates the mathematical coverage argument, the completed search, and independent verification of the emitted tensors. All duality types are included, and no categorifiability filter is applied. Code, machine-readable counts, full multiplication tensors and provenance records accompany the text.
\end{abstract}\vspace{0.4cm}\noindent\textbf{AI assistance and verification.} GPT-6 Astra Pro assisted with code, exposition and computational checks. Exhaustive-search arguments and exact independent tensor verification are recorded separately; no proof-assistant certification or peer-reviewed acceptance is claimed.\vfill\noindent\textbf{Scope.} A census of based rings, not a classification of tensor categories. Only completed whole-rank results are released.\newpage\tableofcontents\newpage

\section{Objects, equivalence and the counting convention}
A \emph{fusion ring} here is a free abelian group with a distinguished finite basis
$B=\{b_0,b_1,\ldots,b_{r-1}\}$, an associative multiplication with unit $b_0$,
and an involution $i\mapsto i^*$ fixing $0$. Its structure constants satisfy
\[
b_i b_j=\sum_k N_{ij}^{k}b_k,\qquad N_{ij}^{k}\in\NN,
\]
\begin{align}
N_{0j}^{k}=N_{j0}^{k}&=\delta_{jk},&N_{ij}^{0}&=\delta_{i^*,j},\\
N_{ij}^{k}&=N_{i^*k}^{j}=N_{kj^*}^{i}.&&\label{eq:reciprocity}
\end{align}
These conventions agree with the usual based-ring definition; see
\cite{EGNO,VS}. No assumption of categorifiability, commutativity or integrality
of Frobenius--Perron dimensions is imposed unless proved in the relevant rank.

A based isomorphism is a bijection of the distinguished bases preserving the
unit and multiplication. In tensor notation it is a permutation $q$ fixing $0$
such that
\[
 N_{ij}^{k}={N'}_{q(i)q(j)}^{q(k)}\quad\hbox{for every }i,j,k.
\]
The coefficient of the unit determines the involution, so such an isomorphism
automatically intertwines duality. The \emph{multiplicity} is the maximum entry
of the \emph{full} ordered multiplication tensor. It is at least one, including
for group rings. Put
\[
 c_r(m)=\#\{\text{rank-}r\text{ fusion rings of multiplicity exactly }m\}/\cong,
 \qquad C_r(M)=\sum_{m=1}^{M}c_r(m).
\]
A list through $M$ contains all rings of multiplicity at most $M$, not only
those of multiplicity equal to $M$. This distinction is essential when exporting
OEIS b-files.

\subsection{Why duality types can be searched separately}
An involution on the $r-1$ nonunit elements consists of $p$ disjoint transpositions
and $r-1-2p$ fixed points, with $0\leq p\leq\lfloor(r-1)/2\rfloor$. Fix one
representative $d_p$ of each type. An isomorphism between two tensors with
involution $d_p$ commutes with $d_p$. Hence it is sufficient to quotient that
stratum by the centralizer
\[
 Z_{S_{r-1}}(d_p),\qquad |Z_{S_{r-1}}(d_p)|=2^p p!(r-1-2p)!.
\]
Different $p$ cannot be isomorphic. Therefore a whole-rank census is complete
exactly when every required stratum has been searched exhaustively. Finishing
only the non-self-dual strata does not certify a whole-rank counting term.
\section{The certified release at rank 3}
The released bound is $M=1000$ and the cumulative number of classes is $C_{3}(1000)=460,353$. Every duality type is included. The accompanying full tensor file has one record per based-isomorphism class.

The first sixteen terms are $4,3,4,6,5,9,6,10,12,9,10,20,9,13,16,25$. The complete 1,000-term file is supplied as \path{oeis/b354471.txt}.

The data and their checksums are recorded in \path{results/census.json}. The corresponding completed-run evidence and independent audit are under \path{verification/rank3/}.

\section{The classical rank-three Diophantine parameterization}
This companion does not claim a new theoretical rank-three classification.
It materializes the familiar parameterization through the requested bound.
In the self-dual basis $(1,X,Y)$, reciprocity gives
\[
X^2=1+mX+kY,\qquad XY=kX+lY,\qquad Y^2=1+lX+nY.
\]
Comparing $(X^2)Y$ with $X(XY)$ proves that associativity is equivalent to
\begin{equation}k^2+l^2=1+lm+kn.\label{eq:rankthree}\end{equation}
Exchanging $X,Y$ sends $(k,l,m,n)$ to $(l,k,n,m)$. Retain $k<l$, or $k=l$
with $m\leq n$.

If $X^*=Y$, reciprocity instead gives
\[
X^2=aX+bY,\quad Y^2=bX+aY,\quad XY=YX=1+aX+aY.
\]
Associativity implies $b^2-a^2=1$. Nonnegative integers force $a=0$, $b=1$,
which is exactly $\ZZ[C_3]$.

\subsection{Every boundary and every arithmetic progression}
For the self-dual part take $k\leq l$. If $k=0$, equation
\eqref{eq:rankthree} gives $l=1$, $m=0$, with $n$ arbitrary in $[0,M]$.
If $k=l$, it gives $k=l=1$ and $m+n=1$, leaving the representative $(1,1,0,1)$.
For $0<k<l$, a common divisor of $k,l$ would divide one, so $\gcd(k,l)=1$.
Set $C=k^2+l^2-1$ and choose $0\leq m_0<k$ with $lm_0\equiv C\pmod k$.
For $k=1$ take $m_0=0$. Then every integral solution is
\[
m=m_0+kt,\qquad n=n_0-lt,\qquad n_0=(C-lm_0)/k,
\]
where
\[
\max\left(0,\left\lceil\frac{n_0-M}{l}\right\rceil\right)
\leq t\leq
\min\left(\left\lfloor\frac{M-m_0}{k}\right\rfloor,
          \left\lfloor\frac{n_0}{l}\right\rfloor\right).
\]
The Euclidean algorithm finds the inverse of $l$ modulo $k$. This visits
every bounded solution exactly once. The two boundary cases and $\ZZ[C_3]$
complete the census. Every emitted solution is checked again in
\eqref{eq:rankthree}; the independent tensor verifier does not use that equation.

\subsection{The deliberate endpoint}
The repository hard-caps the rank-three generator at $M=1000$. It contains
all exact-multiplicity counts and all multiplication tensors through that
bound. There are no rank-three counts or searches beyond that endpoint in
the released data. The b-file belongs to OEIS A354471.

\section{What the independent verification establishes}
The enumerator and verifier have separate mathematical tasks. Exhaustiveness
uses the complete coordinate parameterization, the safe-pruning proof and the
normal completion of every required branch. The verifier reads only the emitted
\emph{full tensors}; it neither reconstructs the parameterization nor trusts the
generator's polynomial list.

For every tensor it checks nonnegative integral coefficients, both unit laws,
uniqueness and involutivity of the dual, both reciprocity identities, and every
integer associativity identity for all four indices. It then canonicalizes
under \emph{every} unit-fixing basis permutation. A matching hash is only a
lookup aid: actual canonical tensors are compared. This detects both literal
repetitions and different labellings of an isomorphic based ring.

Passing these checks proves soundness and uniqueness of the emitted data; by
itself it would not prove that no ring was omitted. That is why normal completion
logs and the mathematical coverage argument are retained separately. Comparison
with the original ancillary file is an additional containment/regression test,
not an input to enumeration and not a replacement for an exhaustiveness proof.

\subsection{Provenance and reproducibility}
The retained Vercleyen--Slingerland v4 ancillary file has SHA-256
\begin{center}\footnotesize\ttfamily
b587c3f683157369da7cc090f762bbff0e18dcf11acc3f5031d6bc0c643843f8
\end{center}
Its 28,451 records contain 25,138 distinct classes. The 3,313 redundant records
occur at rank five and exact multiplicities 9--12. Their original and corrected
counts are respectively
\[
(1463,1794,2283,3049)\quad\longrightarrow\quad(863,1082,1383,1948).
\]
Every deletion has an explicit basis-permutation witness in the audit folder.
This audit does not establish completeness of the source's partial-search cells.

The primary manifest \path{results/census.json} records every released count,
bound, filename and checksum. Compressed tensors use the original line-per-ring,
nested-curly-brace format. An entry is $N[i][j][k]=N_{ij}^{k}$: the output index
has already been dualized where the internal parameter convention requires it.
The raw tensor needs no additional index conversion.

\subsection{Current evidence and release gates}
There is one current dataset per rank under \path{results/rankR/} and one
corresponding evidence directory under \path{verification/rankR/}. Ranks three
and four retain their complete original runs and independent audits. The supplied
laptop censuses at ranks five through eight passed a fresh independent GitHub
audit, run \href{https://github.com/sebastienpalcoux/fusion-ring-census/actions/runs/35553424461}{35553424461}.
Every emitted tensor was checked under all unit-fixing permutations, with zero
duplicate classes and exact agreement with the previous count prefixes. Archive
checksums, source identities, commands, environments and nonempty subprocess
logs are retained. The rank-six upload reproduces the existing bound-seven
census exactly; it is not an additional extension.

The hosted frontier workflow is manual and selects ranks five through eight.
Each rank has one 4,200-second shared driver deadline and a 75-minute job ceiling,
at most 300 runner-minutes in total. Compilation, every attempted multiplicity,
all duality types, export, compression and independent verification share that
deadline. Enumeration receives 90\% of remaining driver time, reserving the
balance for the other phases. Threads follow the actual runner CPU allocation.
The local launcher has no clock limit; its mathematical scope limits and mandatory
verification remain in force.

Only a complete, independently verified whole-rank result can become a candidate.
A subsequent timeout preserves the last verified candidate. Integration requires
explicit review and a fresh independent audit. Interrupted output cannot supply
exhaustive census or OEIS terms. All workflows are manual; a push starts no search.
Earlier payloads and campaigns remain in Git history, without duplicate releases
in the current tree. See \path{docs/PROVENANCE.md} and \path{verification/README.md}.

\section{Reading and reusing the data}
The result files count based rings, not monoidal categories or realizations of
a ring. One tensor can admit several inequivalent categorifications, or none.
No number here is a classification of fusion categories. The rank-specific
count files use exact multiplicity; cumulative totals belong only in the
summary column. Proposed OEIS additions must be regenerated from a complete
manifest, not copied from an interrupted log.

The retained code was generated with AI assistance and is accompanied by explicit
mathematical coverage arguments and independent integer checks. These checks
are not a formal proof in a proof assistant. The present text is an explanatory
computational companion, not a claim of journal acceptance or external peer review.

\section{Reproduction commands}
From the repository root on a local computer, request a complete run with no clock limit. An interruption does not certify a smaller census.
\par\noindent\begin{minipage}{\linewidth}
\begin{lstlisting}
python3 scripts/run_laptop.py --rank 3 --bound 1000
\end{lstlisting}
\end{minipage}\par
To verify the supplied release without recomputing it:
\par\noindent\begin{minipage}{\linewidth}
\begin{lstlisting}
python3 scripts/check_release.py --full
\end{lstlisting}
\end{minipage}\par
For a shared-budget frontier search, use \path{scripts/extend_census.py}. For ranks five through eight it starts at the released bound plus one and restarts the unfinished bound without claiming checkpoint resumption. The local command above verifies all tensors and the old count prefix before packaging a result. The separate timed frontier wrapper bounds every phase. Rank three is instead capped at 1000.

\begin{thebibliography}{9}\small\setlength{\itemsep}{2pt}
\bibitem{EGNO} P. Etingof, S. Gelaki, D. Nikshych and V. Ostrik,
\emph{Tensor Categories}, Mathematical Surveys and Monographs 205,
American Mathematical Society, 2015.
\bibitem{VS} G. Vercleyen and J. Slingerland,
\emph{On Low Rank Fusion Rings}, Journal of Mathematical Physics 64 (2023),
091703; \href{https://arxiv.org/abs/2205.15637v4}{arXiv:2205.15637v4}.
\bibitem{Ostrik} V. Ostrik, \emph{On formal codegrees of fusion categories},
Mathematical Research Letters 16 (2009), 895--901;
\href{https://arxiv.org/abs/0810.3242}{arXiv:0810.3242}.
\bibitem{DP} J. Dong and S. Palcoux, rank-four and rank-five census derivations,
project manuscripts and computational supplements, September 2026.
The relevant derivations are retained in the repository's \path{methods/} directory.
\bibitem{OEIS} OEIS Foundation Inc., \emph{The On-Line Encyclopedia of Integer
Sequences}, entries A348305, A354471, A354472, A354473, A354475, A354476
and A354477. Proposed updates in this repository are not submitted edits.
\end{thebibliography}
\end{document}
